\begin{proof}[Proof of \cref{lem:p213-triv-interval-current}]
Since the ambient dimension is one, a current of bidegree $(0,0)$ has
bidimension $(1,1)$ and acts on compactly supported top-degree forms.
Fix an ordinary affine coordinate \(r\) on a relatively compact open
subinterval \(J\Subset I\) which avoids an infinite endpoint.  A positive
order-zero current on \(J\) is represented by a positive Radon measure
\(\mu\):
\begin{equation}
 \left\langle T,\varphi\,\mathrm d'r\wedge\mathrm d''r\right\rangle
       =\int_J\varphi\,\mathrm d\mu .
\label{eq:p213-triv-interval-measure}
\end{equation}
Here order zero follows directly from positivity; it is also the
one-dimensional instance of \cite[Proposition~7.2.9]{CLD25}.

For every \(q\in C_c^\infty(J)\), closedness gives
\begin{equation}
 0=\left\langle T,
       \mathrm d'\bigl(q\,\mathrm d''r\bigr)\right\rangle
   =\int_Jq'(r)\,\mathrm d\mu(r).
\label{eq:p213-triv-zero-distributional-derivative}
\end{equation}
Thus the distributional derivative of \(\mu\) is zero.  A distribution on a
connected interval with zero first derivative is constant, so
\(\mu=c\,\mathrm dr\) for a constant \(c\geq0\).  Constants obtained on
overlapping subintervals agree.  This proves
\eqref{eq:p213-triv-interval-classification}.

At the infinite endpoint of the modified interval in
\cite[\S6.7.2]{CLD25}, all paralinear coordinates are constant near the
endpoint.  Top-degree forms therefore vanish near that endpoint, and no
additional endpoint current occurs.  The same observation applies at every
included finite endpoint in the modified structure: paralinear functions
are constant there by definition, so positive-degree test forms vanish
near it.  In particular, no Stokes boundary term is visible there.  We do
not use or assert a zero-current conclusion for an interval with tropical
boundary.

Finally, \cite[Example~8.6.4]{CLD25} states that the integration current is
Bedford--Taylor.  The cone of Bedford--Taylor currents is stable under
nonnegative scalar multiplication, proving the last assertion.
\end{proof}

\medskip

We use the following elementary local computation twice.  It is included
to prevent a hidden skeletal-current assumption.

\begin{lemma}
\label{lem:p213-triv-fan-normal-form}
Let
\[
 \Sigma=\bigcup_{j=1}^m\mathbf R_{\geq0}v_j\subset\mathbf R^N
\]
be a finite one-dimensional fan with pairwise distinct rays.  Let \(S\) be
an order-zero positive supercurrent of bidimension \((1,1)\), supported on
\(\Sigma\), and assume \(\mathrm d'S=\mathrm d''S=0\).  Then the restriction
of \(S\) to the \(j\)-th open ray is
\begin{equation}
              w_j\,[\mathbf R_{>0}v_j],\qquad w_j\geq0,
\label{eq:p213-triv-ray-density}
\end{equation}
there is no current component supported at the vertex, and
\begin{equation}
                         \sum_{j=1}^m w_jv_j=0.
\label{eq:p213-triv-fan-balance}
\end{equation}
\end{lemma}

\begin{proof}
On the relative interior of a ray, choose linear coordinates
\(y_1,\ldots,y_N\) with \(y_1\) increasing in the primitive ray direction
and \(y_2=\cdots=y_N=0\) on the ray.  An order-zero positive supercurrent of
bidimension \((1,1)\) is represented there by a positive-semidefinite
matrix of Radon measures.  Test \(\mathrm d'S=0\) against forms whose
coefficients have arbitrary first jets in the normal directions.  Every
row and column except the \(y_1y_1\)-entry vanishes.  The remaining equation
is the distributional identity
\[
 \partial_{y_1}\mu=0,
\]
hence \(\mu=w_j\,\mathrm dy_1\) with \(w_j\geq0\).  This proves
\eqref{eq:p213-triv-ray-density} without assuming that \(S\) was skeletal.

After removing these ray parts, an order-zero positive remainder at the
vertex would have the form \(A\delta_0\), with \(A\) a positive-semidefinite
matrix.  Its differential contains first derivatives of \(\delta_0\),
whereas the boundary of each ray current contains only a multiple of
\(\delta_0\).  Tests with arbitrary first jet at the origin force \(A=0\).
Testing closedness against one-forms that are constant in a neighborhood
of the vertex now gives exactly
\eqref{eq:p213-triv-fan-balance}.  The \(\mathrm d''\)-equation is equivalent
by symmetry.  This is the one-dimensional local calculation behind the
harmonicity condition of \cite[Theorem~6.6.8 and
Remark~7.7.11]{CLD25}.
\end{proof}

\begin{proof}[Proof of
\cref{prop:p213-triv-abstract-cross-counterexample}]
The coordinate cross is a one-dimensional tropical space without boundary:
at the origin its four primitive outgoing vectors satisfy
\begin{equation}
                  e_1-e_1+e_2-e_2=0.
\label{eq:p213-triv-unit-cross-balance}
\end{equation}
The current in \eqref{eq:p213-triv-abstract-cross-current} is positive by
construction.  On the relative interior of each ray its density is constant.
At the origin its weighted outgoing-vector sum is
\begin{equation}
                  2e_1-2e_1+e_2-e_2=0.
\label{eq:p213-triv-weighted-cross-balance}
\end{equation}
Stokes' formula on the four rays therefore gives
\(\mathrm d'T=\mathrm d''T=0\); equivalently, this is the sufficient
direction of the balancing computation in
\cref{lem:p213-triv-fan-normal-form}.

The inclusion \(\Sigma\hookrightarrow\mathbf R^2\) is a tropical chart, and
the two ambient functions in
\eqref{eq:p213-triv-abstract-cross-potentials} are continuous and convex.
Their restrictions are psh by
\cite[Proposition~8.2.5]{CLD25}.  For a continuous piecewise-affine function
\(q\) on a balanced one-dimensional fan, integration by parts gives
\begin{equation}
 \mathrm d'\mathrm d''[q]
   =\left(\sum_{\rho}\partial_{\rho}q(0)\right)\delta_0,
\label{eq:p213-triv-unweighted-corner-formula}
\end{equation}
when \(q\) is affine away from the vertex; the sign is the one for which a
convex corner is positive.  The outgoing slope vectors of \(U\) and \(V\),
in the order
\(\rho_{+1},\rho_{-1},\rho_{+2},\rho_{-2}\), are respectively
\begin{equation}
                         (2,2,0,0),
       \qquad            (1,1,1,1).
\label{eq:p213-triv-abstract-cross-slopes}
\end{equation}
Their sums are both \(4\), proving the first equality in
\eqref{eq:p213-triv-abstract-cross-defect}.

Multiplication by \(T\) replaces the unit ray densities in
\eqref{eq:p213-triv-unweighted-corner-formula} by \(2,2,1,1\).  Hence
\begin{equation}
 \begin{aligned}
  \mathrm d'\mathrm d''(UT)
     &=(2\cdot2+2\cdot2+1\cdot0+1\cdot0)\delta_0
       =8\delta_0,\\
  \mathrm d'\mathrm d''(VT)
     &=(2\cdot1+2\cdot1+1\cdot1+1\cdot1)\delta_0
       =6\delta_0.
 \end{aligned}
\label{eq:p213-triv-weighted-corner-computation}
\end{equation}
This violates \cite[Definition~8.6.2]{CLD25}.  Notice that the underlying
topological space is a four-ended \(\mathbf R\)-tree; the failure comes from
the two-dimensional cone of nonnegative balanced calibrations at its central
vertex, not from a cycle.
\end{proof}

\begin{proof}[Proof of \cref{prop:p213-triv-p1-current-rigidity}]
Chambert--Loir and Ducros give an exact model of \(X\) in
\cite[\S6.7.3]{CLD25}.  Write \(B=\mathbf P^1(k)\).  The underlying space is
the quotient of
\[
 B\times[-\infty,0]
\]
obtained by identifying all \((b,0)\) with a central point \(\eta\).  Let
\(X_b\) be the branch indexed by \(b\).  A smooth function is smooth on
each branch, is constant near a rigid endpoint, is nonconstant on only
finitely many full branches near \(\eta\), and its outgoing derivatives at
\(\eta\) have finite support and sum zero.  The resulting tropical space is
boundaryless and is realized by
\((\mathbf P^1_k)^{\mathrm{an}}\).

First suppose that a connected open \(Y\) does not contain \(\eta\).  It is
contained in a single branch and
\cref{lem:p213-triv-interval-current} applies.

Suppose now that \(\eta\in Y\).  On the open interior of \(X_b\cap Y\),
\cref{lem:p213-triv-interval-current} gives a constant branch density
\(c_b\geq0\).  We show that these constants are all equal.  Fix distinct
indices
\[
                    b_0,b_1,\ldots,b_m\in B.
\]
The topology in \cite[\S6.7.3]{CLD25} says that \(Y\) contains every full
branch except possibly finitely many, and contains an initial segment of
each exceptional branch.  Thus we may choose a compact tropical domain
\(W\subset Y\) containing \(\eta\) and initial segments of the selected
branches.  The cutoff needed to localize \(T\) is explicit: take it equal to
\(1\) on every nonexceptional full branch, and on each of the finitely many
exceptional branches take a one-variable smooth cutoff which equals \(1\)
near \(\eta\) and vanishes before the omitted tail.  The criterion in
\cite[\S6.7.3]{CLD25} makes this a smooth function.  Its support is closed in
the compact star and contained in \(Y\), hence compact.  Thus the localized
chart pushforwards used below are properly supported.

Put \(r=-t\), so that \(r\) is the outgoing coordinate on a branch.  For
\(1\leq i\leq m\), choose the rational function
\[
                     F_i\in k(\mathbf P^1)^\times,
       \qquad        \operatorname{div}(F_i)=b_i-b_0,
\]
and shrink only the selected branch segments so that \(F_i\) is invertible
on \(W\).  Then
\[
                         \ell_i=-\log|F_i|
\]
has outgoing slope \(1\) on \(b_i\), slope \(-1\) on \(b_0\), and slope
zero on every other branch.  Hence these are actual analytic moment
coordinates, not merely affine functions on a coarse abstract star.  Their
slope sum is zero, in agreement with the explicit affine-function
description in \cite[\S6.7.3]{CLD25}.

The chart
\[
             \ell=(\ell_1,\ldots,\ell_m)
\]
has, at $\ell(\eta)=0$, ray vectors
\begin{equation}
       v_{b_i}=e_i\quad(1\leq i\leq m),
       \qquad
       v_{b_0}=-(e_1+\cdots+e_m).
\label{eq:p213-triv-simplex-rays}
\end{equation}
Every other branch is contracted by this chart.  Pullback of a top-degree
form through a constant branch is zero, so those branches contribute
nothing to the pushforward; in particular, they do not become a spurious
Dirac current at the vertex.  Apply
\cref{lem:p213-triv-fan-normal-form} to the localized chart representation
of \(T\).
Its balancing equation is
\begin{equation}
       \sum_{i=1}^mc_{b_i}e_i
          -c_{b_0}(e_1+\cdots+e_m)=0.
\label{eq:p213-triv-simplex-balance}
\end{equation}
Thus \(c_{b_i}=c_{b_0}\) for every \(i\).  Since the finite set of indices
was arbitrary, all branch densities are a single constant \(c\geq0\).  No
summability or boundedness of the a priori family \((c_b)_{b\in B}\) was
used: each tropical chart and each top-degree test form sees only finitely
many noncontracted branches.

The same local normal-form lemma excludes an additional order-zero current
supported at \(\eta\).  Here is the chart-level point.  Every top-degree test
form near \(\eta\) is represented on a finite analytic moment chart.  In
that chart the atomic part of the localized pushforward is a
positive-semidefinite matrix times \(\delta_{\ell(\eta)}\); the first-jet
test in \cref{lem:p213-triv-fan-normal-form} makes this matrix zero.  A
branch contracted by the chart cannot create such an atom, since the
pullback of every top-degree form to a constant branch is zero.  Running
this argument in the particular finite chart representing the test form
shows that a current supported at \(\eta\) evaluates to zero on every test
form.

Top-degree forms vanish near the rigid endpoints, because all tropical
coordinates are constant there.  Hence there is no endpoint contribution
either, and \(T=c\delta_Y\).  This proves
\eqref{eq:p213-triv-p1-classification}.

Open subsets of a boundaryless tropical space are boundaryless.  The
restriction of $\delta_Y$ to an open subset is its integration current.
The full Bedford--Taylor assertion therefore follows again from
\cite[Example~8.6.4]{CLD25}.  The descriptions of the affine line and the
disc and annulus pieces are precisely the realizations in
\cite[\S\S6.7.1--6.7.3]{CLD25}.
\end{proof}

\begin{proof}[Proof of
\cref{prop:p213-triv-algebraic-curve-current-rigidity}]
Let \(\eta_C\) denote the trivial absolute value on \(k(C)\).  Every
nontrivial real valuation of \(k(C)\) which is trivial on \(k\) has a center
at a closed point \(p\in C\).  Since \(C\) is regular of dimension one,
\(\mathcal O_{C,p}\) is a discrete valuation ring; the valuation is a
positive real multiple of \(\operatorname{ord}_p\).  Thus
\(C^{\mathrm{an}}\) is a star with center \(\eta_C\), one ray for every
closed point \(p\), and a rigid endpoint at the end of that ray.  This is
also the type--\(2\) residue-curve description of
\cite[\S3.3.5.2, Theorem~4.2.10 and \S4.2.11.1]
{Ducros24CurveStructure}: at an interior type--\(2\) point, branches are in
bijection with the closed points of the residue curve.

As in the interval proof, positivity and closedness make the restriction
of \(T\) to the ray at \(p\) a constant multiple \(w_p\,\mathrm dr\), with
\(w_p\geq0\).  If \(f\in k(C)^\times\), then the tropical coordinate
\(-\log|f|\) has outgoing slope \(\operatorname{ord}_p(f)\) on the
\(p\)-ray, up to one common harmless sign convention.  This is also the
one-dimensional divisor-slope instance of the Poincar\'e--Lelong formula
\cite[Theorem~16.7.3]{CLD25}.  Testing closedness
at \(\eta_C\) gives
\begin{equation}
                       \sum_{p\in C(k)}
                       w_p\operatorname{ord}_p(f)=0.
\label{eq:p213-triv-principal-annihilation}
\end{equation}
The sum is finite.  Equivalently, the positive linear functional
\begin{equation}
      L\colon\operatorname{Div}(C)\longrightarrow\mathbf R,
      \qquad
      L\!\left(\sum_pn_pp\right)=\sum_pn_pw_p
\label{eq:p213-triv-divisor-functional}
\end{equation}
annihilates principal divisors and descends to \(\operatorname{Pic}(C)\).

We now use positivity, not a claim that every degree-zero divisor is
principal.  Fix \(p_0\in C(k)\), write \(g\) for the genus, and choose a
single integer \(N\geq2g-1\).  For every degree-zero divisor \(D\),
Riemann--Roch gives effective divisors \(E,E'\) such that
\begin{equation}
                     D+Np_0\sim E,
       \qquad       -D+Np_0\sim E'.
\label{eq:p213-triv-rr-effective}
\end{equation}
See \cite[Tag 0B90]{StacksProject}.  Since \(L\) is nonnegative on effective
divisors,
\begin{equation}
                         |L(D)|\leq Nw_{p_0}.
\label{eq:p213-triv-pic0-bound}
\end{equation}
Apply the same estimate to \(mD\), for every positive integer \(m\).  It
follows that
\[
 |mL(D)|\leq Nw_{p_0}
\]
for all \(m\), hence \(L(D)=0\).  Taking \(D=p-p_0\) proves
\begin{equation}
                              w_p=w_{p_0}
                 \qquad\text{for every }p\in C(k).
\label{eq:p213-triv-all-curve-weights}
\end{equation}

For completeness, this argument is insensitive to the genus: the point is
that Riemann--Roch supplies a uniform effective translate of every
degree-zero class.  It also shows why a nonzero homomorphism
\(\operatorname{Pic}^0(C)\to\mathbf R\) cannot furnish a positive current;
such a homomorphism is necessarily unbounded on its integer multiples,
whereas positivity and \eqref{eq:p213-triv-rr-effective} give the uniform
bound \eqref{eq:p213-triv-pic0-bound}.

The finite-chart argument in
\cref{lem:p213-triv-fan-normal-form} rules out a residual current at
\(\eta_C\).  More explicitly, take the actual finite analytic moment chart
on which a given top-degree test form is represented.  Only finitely many
rays are noncontracted in that chart.  The local fan calculation kills its
vertex matrix, while a contracted branch has zero pullback on top-degree
forms and hence contributes no artificial vertex atom.  Since every test
form is represented on an actual finite chart, this proves vanishing of the
putative point-supported current without asserting that one fixed chart
sees all branches.  Branches can moreover be separated by analytic moment
functions by
\cite[Lemma~3.3.6 and Theorem~4.2.5]{Ducros24CurveStructure}.
Top-degree forms vanish near the rigid endpoints.  Therefore
\(T=c\delta_{C^{\mathrm{an}}}\), where \(c=w_{p_0}\), proving
\eqref{eq:p213-triv-curve-classification}.  The Bedford--Taylor conclusion
follows from \cite[Example~8.6.4]{CLD25}.

Removing finitely many closed points merely replaces the corresponding
closed rays by open rays.  The calculation at \(\eta_C\), including the
Riemann--Roch bound, is unchanged, and compactly supported tests introduce
no term at a missing endpoint.  This proves the Zariski-open assertion.
We do not use this sentence to claim arbitrary analytic opens of
\(C^{\mathrm{an}}\): for positive genus that stronger localization requires
a separate audit of the induced analytic affine structure on an arbitrary
star-shaped neighborhood.
\end{proof}

\begin{proof}[Proof of \cref{prop:p213-triv-four-puncture-no-go}]
The explicit star description of \cite[\S6.7.3]{CLD25} gives the slope
table
\begin{equation}
\begin{array}{c|rrrr}
 &a&b&c&d\\ \hline
 s=-\log|f|&1&-1&0&0\\
 t=-\log|g|&0&0&1&-1.
\end{array}
\label{eq:p213-triv-cross-slopes}
\end{equation}
All other branches are contracted.  Hence the tropical image is the full
cross with its four unit weights.  The ambient functions \(2|s|\) and
\(|s|+|t|\) are continuous convex functions.  Their pullbacks are psh by
\cite[Proposition~8.2.5]{CLD25}; equivalently one may use the decreasing
smooth convex approximants
\begin{equation}
 2\sqrt{s^2+\varepsilon},
 \qquad
 \sqrt{s^2+\varepsilon}+\sqrt{t^2+\varepsilon}
 \quad(\varepsilon\downarrow0).
\label{eq:p213-triv-cross-approximants}
\end{equation}
The outgoing slopes at \(\eta\) are
\begin{equation}
\begin{array}{c|rrrr|c}
 &a&b&c&d&\text{sum}\\ \hline
 U&2&2&0&0&4\\
 V&1&1&1&1&4.
\end{array}
\label{eq:p213-triv-cross-potential-slopes}
\end{equation}
There is no curvature on an open ray and no contribution on a contracted
branch, proving \eqref{eq:p213-triv-cross-curvature}.

The coarse vectors admit the apparently balanced weights
\((2,2,1,1)\).  Refine the moment map by
\[
                         h=\frac{z-a}{z-c}.
\]
The ray vectors of
\[
 (-\log|f|,-\log|g|,-\log|h|)
\]
are
\begin{equation}
 v_a=(1,0,1),\quad v_b=(-1,0,0),\quad
 v_c=(0,1,-1),\quad v_d=(0,-1,0).
\label{eq:p213-triv-refined-rays}
\end{equation}
If the visible weights of a positive closed current are
\(w_a,w_b,w_c,w_d\), the local normal form
\cref{lem:p213-triv-fan-normal-form} gives
\begin{equation}
                       \sum_{e=a,b,c,d}w_ev_e=0.
\label{eq:p213-triv-refined-balance}
\end{equation}
In coordinates, this is
\begin{equation}
                    (w_a-w_b,\;w_c-w_d,\;w_a-w_c)=0,
\label{eq:p213-triv-refined-equalities}
\end{equation}
so all four weights are equal.  The proposed tuple \((2,2,1,1)\) has
refined defect \((0,0,1)\).

One can see the same obstruction directly on the star.  Let \(T_w\) be the
positive branch-weight current having weights \((2,2,1,1)\) on
\(a,b,c,d\) and any bounded nonnegative weights on the contracted branches.
Put \(H=U-V\).  Its four slopes are \((1,1,-1,-1)\), so the weighted corner
calculation gives
\begin{equation}
                   \mathrm d'\mathrm d''(H T_w)
                       =(2+2-1-1)\delta_\eta
                       =2\delta_\eta.
\label{eq:p213-triv-visible-bt-defect}
\end{equation}
But the same refined imbalance appears already in
\(\mathrm d'T_w\) and \(\mathrm d''T_w\): \(T_w\) is not closed.  Thus
\eqref{eq:p213-triv-visible-bt-defect} is not a Bedford--Taylor
counterexample.  It is the diagnostic showing exactly why the natural
candidate is inadmissible.
\end{proof}
